\subsection{Scaling Graph Neural Networks to Large Graphs}
\label{sec:relwork:scaling}

\subsubsection{Sampling-Based Approaches}
\label{sec:relwork:scaling:sampling}
The oldest answer to a graph that does not fit is to sample it: neighbour sampling caps
the fanout of each aggregation \citep{hamilton2017graphsage}, and subgraph sampling
trains on induced pieces. Both are per-epoch stochastic operations built for node-level
targets, and neither survives the move to this task: graph-level regression needs every
node in one forward pass (\ref{sec:prelim:gnn:readout}), and the reduction under study
must be a deterministic, cacheable artifact computed once, not redrawn every epoch.

\subsubsection{Partitioning and Distributed Training}
\label{sec:relwork:scaling:distributed}
The standard industrial answer is to partition the graph across devices and exchange
messages at the boundaries, or drop them. Cutting edges is then forced rather than
chosen, which is why partitioning is carried as a reduction family here
(\ref{sec:prelim:reduction:partitioning}) despite removing no nodes: the accuracy cost of
a severed spanning edge is exactly what a distributed practitioner pays, measured on one
device. DistDGL \citep{zheng2020distdgl} represents the practice at billion-node scale.

\subsubsection{Memory-Efficient Training Techniques}
\label{sec:relwork:scaling:memory}
Gradient checkpointing, mixed precision and gradient accumulation, defined
in~\ref{sec:prelim:gnn:cost}, are used here rather than studied here, but a reader will
still ask why reduction is needed when checkpointing exists. The answer is composition.
Each technique trades compute, precision or latency for memory at a fixed input size;
reduction changes the input size itself, and the two stack. The baseline configurations
depend on the first fact, DeepGate4 being runnable only under checkpointing
(\ref{sec:method:architecture:baselines}), and the contribution depends on the second.